\documentclass[11pt]{article}
\usepackage{mathblatt}
\begin{document}
\blattkopf*{Basisaufgaben}{BAS-Z10}{Basisaufgaben · Zettel · BAS-Z10}
\einheitenkopf[][Zettel 10]{Basisaufgaben · Zettel 10}
\zweigzeile{ohne Taschenrechner · je Aufgabe 1 Punkt · Lösungen auf der Rückseite}
% \small: zehn Aufgaben mit bis zu zwei Koordinatensystemen auf einer Seite (Render 28.09.)
\small

% 1: Bruchteil einer Fläche bestimmen – brueche-dezimalzahlen-basis-k1-v9 (Original 2026-FOR-B1b)
\begin{aufgabe}{Bei welcher Figur ist genau $\frac{3}{10}$ grau? \hfill (P26F) \\ \kreuz{P} \kreuz{Q} \kreuz{R} \kreuz{S}}

P \bruchkreis[0.8]{3}{8} \quad Q \kreissektor[0.8]{108}{} \quad R \bruchrechteck[2.5]{3}{9} \quad S \kreissektor[0.8]{30}{}
\end{aufgabe}

% 2: Zahlen in verschiedenen Darstellungen vergleichen – brueche-dezimalzahlen-basis-k4-v8 (Original 2023-OS-B1f)
\begin{aufgabe}{Welcher Wert ist der kleinste: $2{,}2$; $0{,}22$; $22\,\%$; $0{,}2^2$? \hfill (P23) \\ \leerfeld}
\end{aufgabe}

% 3: Termwert berechnen – rationale-zahlen-basis-k1-v6 (Original 2026-FOR-B1g)
\begin{aufgabe}{Berechne den Wert von $5 \cdot (x + 2)$ für $x = -5$. \hfill (P26F) \\ \leerfeld}
\end{aufgabe}

% 4: Exponent einer Potenz bestimmen – potenzen-wurzeln-basis-k1-v3 (Original 2024-OS-B1g)
\begin{aufgabe}{$5^x = 25$ – welche Zahl ist $x$? \hfill (P24) \\ x = \leerfeld}
\end{aufgabe}

% 5: Term zu Sachtext angeben – terme-basis-k2-v5 (Original 2023-OS-B1h)
\begin{aufgabe}{Die Summe aus einer Zahl $n$ und 8 wird mit 5 multipliziert. Welcher Term passt? \hfill (P23) \\ \kreuz{$5n + 8$} \kreuz{$5 \cdot (n + 8)$} \kreuz{$n + 8 \cdot 5$} \kreuz{$(n + 5) \cdot 8$}}
\end{aufgabe}

% 6: Lösung durch Einsetzen prüfen – quadratische-gleichungen-basis-k1-v4 (Original 2025-OS-B1h)
\begin{aufgabe}{Welcher Wert erfüllt $x \cdot (x + 9) = -20$? \hfill (P25) \\ \kreuz{$x = 4$} \kreuz{$x = 5$} \kreuz{$x = -5$} \kreuz{$x = -20$}}
\end{aufgabe}

% 7: y-Achsenabschnitt ablesen – lineare-funktionen-basis-k4-v1 (Original 2024-OS-B1i)
\begin{aufgabe}{\begin{minipage}[t]{0.6\linewidth}Schnittpunkt mit der y-Achse: Markiere ihn an der Geraden $y = 2x + 1$. \hfill (P24)\end{minipage}\hfill\begin{minipage}[t]{0.37\linewidth}\vspace{0pt}
\begin{ksys}[xmin=-5,xmax=5,ymin=-5,ymax=5,klein] \gerade{2}{1}{g} \end{ksys}
\end{minipage}}
\end{aufgabe}

% 8: Figur nach Spiegelung benennen – symmetrie-abbildungen-basis-k1-v2 (Original 2018-OS-B1h)
\begin{aufgabe}{Ein rechtwinkliges Dreieck liegt mit einer Kathete auf der Geraden $g$ und wird an $g$ gespiegelt. Welche Figur bilden Dreieck und Spiegelbild? \hfill (P18)}
\end{aufgabe}

% 9: Dreiecksungleichung anwenden – winkel-dreiecke-basis-k1-v2 (Original 2020-OS-B1i)
\begin{aufgabe}{Aus welchen drei Stäben kann man ein Dreieck legen? \hfill (P20) \\ \kreuz{$2$ cm, $3$ cm, $6$ cm} \kreuz{$4$ cm, $5$ cm, $8$ cm} \kreuz{$3$ cm, $4$ cm, $7$ cm}}
\end{aufgabe}

% 10: Gegenfläche im Würfelnetz bestimmen – koerper-basis-k1-v2 (Original 2016-OS-B1j)
\begin{aufgabe}{\begin{minipage}[t]{0.6\linewidth}Das Netz aus sechs Quadraten (je Buchstabe ein Quadrat) wird zu einem Würfel gefaltet. Die Fläche T ist der Boden einer Schachtel. Welche Fläche liegt oben? \hfill (P16) \\ \leerfeld\end{minipage}\hfill\begin{minipage}[t]{0.37\linewidth}\vspace{0pt}
$\begin{array}{cccc}  & P &  &  \\ Q & R & S & T \\  &  & U &  \end{array}$
\end{minipage}}
\end{aufgabe}

\begleitteil
\erg{1}{Q, denn $\frac{108}{360} = \frac{3}{10}$}
\erg{2}{$0{,}2^2 = 0{,}04$; die anderen sind $2{,}2$ und zweimal $0{,}22$}
\erg{3}{$5 \cdot (-5 + 2) = 5 \cdot (-3) = -15$}
\erg{4}{$5, 25$: zwei Faktoren, also $x = 2$}
\erg{5}{$5 \cdot (n + 8)$}
\erg{6}{$x = -5$, denn $(-5) \cdot (-5 + 9) = (-5) \cdot 4 = -20$}
\erg{7}{$(0|1)$}
\erg{8}{gleichschenkliges Dreieck}
\erg{9}{$4$ cm, $5$ cm, $8$ cm}
\erg{10}{R (R und T liegen sich gegenüber)}
\end{document}
